\subsection{Additional Native Pruning and Relevance Results}\label{sec:newpruning}
We distinguish the quality of the selected native representation from
the quality obtained by transferring its dimension to an ordinary VAE.
Counts alone do not measure successful compression. Historical versus
additional readouts are collected in Appendix~\ref{app:gpu}.

Table~\ref{tab:armnew} separates native reconstruction, dimension transfer
and endpoint constraint satisfaction. For MNIST and Fashion-MNIST,
all three native posterior-mean reconstructions meet $Q$, while only
two training constraints do. All three dSprites native runs fail both
checks, although their transferred ordinary VAEs meet $Q$. Their native
MSE is approximately 0.04236, and all 64 gate probabilities remain
between 0.1 and 0.9 (Appendix~\ref{app:gpu}); thresholded counts of 48--62
therefore do not demonstrate settled sparse solutions. CIFAR-10 meets
native mean-reconstruction $Q$ in one seed and the training constraint
in two. These are limitations of the local implementation.
@@table52:arm_native@@

The sampled native ARM MSE exceeds the deterministic anchor-derived
target in every run. This comparison demonstrates the consequence of
changing the reconstruction definition; it is not a selection study
under a target derived from a sampled anchor. It also explains why
training-constraint status cannot be inferred from mean-reconstruction $Q$.

The exact-Jacobian ARD zero-centre readouts are 6.0, 4.0, 9.0 and 26.7
on MNIST, Fashion-MNIST, dSprites and CIFAR-10, respectively
(Table~\ref{tab:ardnew}). The sample-centred variant gives 6.0, 4.0, 9.7
and 26.3. Small differences across these three seeds support a bounded
centering-sensitivity observation, not invariance to prior specification.
Both settings meet full-model and transfer $Q$ on the three datasets
other than dSprites and fail both on dSprites. These initial evaluations
do not mask the selected-out axes.
@@table52:ard_native@@

\paragraph{Quality after retaining the selected ARD axes.}
Table~\ref{tab:ardpruned} closes that evaluation gap for the
sample-centred variant. Masking changes mean validation MSE by
$+0.045\%$, $+0.017\%$, $+126.3\%$ and $+3.4\%$ on MNIST,
Fashion-MNIST, dSprites and CIFAR-10, respectively. These percentages
are ratios of three-seed mean MSEs relative to the corresponding full
model, not seed-averaged percentage changes. The full, masked, adapted
and transferred models all meet $Q$ in all three seeds on each dataset
other than dSprites; all four fail $Q$ in every dSprites seed.
@@table53:native_pruning_validation@@

For dSprites, masking raises MSE from 0.00133 to 0.00301; decoder
adaptation lowers it to 0.00183, still 37.1\% above the full model.
Thus the local 99\% relevance cutoff does not ensure reconstruction
preservation after jointly removing the other axes. This result does
not identify the contribution of each removed axis and does not
establish failure of the published ARD method. The transferred
ordinary-VAE MSE is 0.00164, which also misses $Q$; its better quality
than the masked model is not evidence that the selected count suffices
for the specified quality target.

For CIFAR-10, the masked and adapted means (0.01489 and 0.01473)
are below the transferred mean (0.01519), although above the full-model
mean (0.01440). MNIST instead favours transfer on this metric.
These descriptive, three-seed results depend on the trained model,
inference rule and adaptation budget; they do not establish general
native-pruning superiority. Appendix~\ref{app:gpu} gives per-seed
values, held-out test reconstruction and partial timing.
